\chapter[Sermon 42. What Should Be Done to the Residue of the Impenitent, in This Mean While Feeling No Evil, of the Locusts and Horses.]{What Should Be Done to the Residue of the Impenitent, in This Mean While Feeling No Evil, of the Locusts and Horses.}
\label{sermon:042}
\markright{Sermon 42. What Should Be Done to the Residue of the Impenitent ...}

And the remnant of the men who were not killed by these plagues did not repent of their actions, so that they would not worship devils and images of gold, silver, bronze, and stone, and of wood—images that cannot see, hear, or go. They also did not repent of their murders, their sorcery, their adultery, or their theft.

\index{Arethas of Caesarea}It is spoken abundantly about how great calamity shall come upon the world of the locusts and horses under the fifteenth and sixteenth trumpets. And where it is sufficiently known that not all are subject to the locusts and horses, nor to be punished by them, even though some commit things worthy of punishment, one might marvel whether those who are free and exempted from these plagues may safely lead an impenitent life. He foresees and says, “And the remainder of men, who also commit shameful things against God, and yet are not slain with these plagues set forth, may not think to escape unpunished.” For even they shall be punished also by God, most just. For the speech is defective, and therefore to be made up, both by the tenor hereof, and also by the catholic sense of the whole Scripture, which is that all impenitent persons are punished by God, and that so much more grievously, the more carelessly they have abused God’s longsuffering, being nothing moved by any examples of God’s judgments. Yet he says not this by express words. It was enough for him to rehearse the wickedness wherein they were drowned. For from this, every man may gather what is due to such offenders. Arethas, a Greek expounder, expounding this place, says this speech shows an excellence of insensibility, that is, of the wantonness and lasciviousness of those who have spent the time granted them by God to repent, in vanity, that even for the worthiness of their slothfulness they might receive their reward; yea, even before the eyes of the ungodly, the very reward is put in effect. Yet these men, not only by the sight of these terrible things which they had present before their eyes, were made never a whit better, but also worse, and more and more wrapped in sin, have fulfilled their course, and so forth.

Hereof we may gather that it is not sufficient for a godly and blessed life that a man be not a Romanist or a Muslim; but that of every one of us is required a true faith, which may make us to walk in all the commandments of God; and that we should know that all must be grievously punished by God, so many as transgress the law of God, of what religion, condition, age, state, or degree so ever they may be. For God, most just, has no respect of persons. Whoever has sinned without a law, says the Apostle, shall perish without law; and whoever has sinned in the law, by the law shall be judged. Certainly John seems here to bring forth both the tables of the law, and thereby to reprove the sins and wickedness of the ungodly men, of whom he will also that judgment be gathered. The first table sets forth the service of God, commanding to worship one God, not to worship idols, and so forth. The second gives precepts of living, and teaches the love of our neighbor, forbidding murder, adultery, theft, and like mischiefs. John brings forth two sins done against the first table, and three or four committed against the second. Neither is there any doubt but that he comprises under these all like or unlike, more or less offenses against God and his will. Whoever therefore you are, if you offend against the divine law, you shall be punished. If you seem in this world to escape freely and to flit from hence happily, the same may chance unto you that happened to the rich glutton, whose judgment is described in Luke 16. Briefly, he shall be punished, whoever offends God. God knows the manner, whether he shall punish here and in the world to come, or in the world to come only, and grant here a [illegible] voluptuous life.

\index{Repentance}And we must chiefly observe in this treatise that sinners are not here condemned. For we are all sinners; otherwise no one would be saved. Damned are those who do not repent, those who truly die in their sins without repentance. The Apostle denies that idolaters, adulterers, thieves, covetous persons, extortioners, and so on, will possess the kingdom of God, but he adds: “But such you were, but you have been washed, but you have been sanctified, but you have been justified in the name of our Lord Jesus, and by the Spirit of our God.” And if you doubt whether you may again come into favor with God, if you, having once been enlightened and justified, fall again into sin, learn from the fall and sin of Saint Peter that you may be restored. And as it is written, “Seven times the righteous falls, and rises again,” and so on. Therefore let us learn from this how effective repentance is, and how harmful a lack of repentance is. If you are, or have been, an idolater, you ought not to despair; turn to the Lord and do penance. If you fall again, do not remain in your wickedness. I have spoken more about this in another place. But if you will not return to God, nor leave the evil custom of sin, never look for any grace of God. You shall perish in your sins.

It remains that we declare in few words the forms of sins, set forth here by John, under which, as I said before, he has undoubtedly included similar offenses, so that the same judgment may be had of like things. First, he says, as it were generally, that they have not repented of the works of their hands. For although with this note or mark idolatry is condemned in the prophets, I extend it to all other deeds proceeding from the force of the flesh. For our work is truly sin; and the good work is of the grace of God and of regeneration. And this general thing once set forth, he adds diverse kinds and forms, two against the first table, and four or three against the second.

It is against the first commandment to worship devils. For our very God will have himself alone taken for God, honored and worshipped. And who is so mad, say you, that will worship devils? Verily, there are certain people in the East, who are said to worship devils, for no other end but that they should not hurt them. This is a barbarous and foolish people; why do they not rather worship him who is only able to restrain the devil, that he can not hurt? However, this wickedness stretches far. For they indeed worship the devil, which will seem to worship God. For this matter is not esteemed of the opinion or intent of the worshipper, but of the law maker. For the gentiles would not seem that they sacrificed to devils, but would have taken it most displeasingly, if any should have said that they worshipped the devil. “You are a most vile and most impudent varlet and slanderer,” they would have said, “which dare so reproach the gods and us.” But Paul nevertheless, I say not, says he, that an idol, or that which is offered unto idols, is anything: but this I say, that the things which the gentiles offer up, they offer them to devils, and not to God. For where there is one only God, and he allows only these sacrifices, which are offered to him, he calls strange gods devils, and idol offerings sacrificed to the devil: of this judgment is the thing esteemed, and not of the fond intent of men. King Saul would have offered to God the burnt offering of Samuel: but Samuel told him, that he should commit idolatry and magic, and so forth. This is a hard saying, but yet true. Whereof I have spoken in another place more at large. The worshipping of images of God and of the saints is against the precept of the first table. For all idolatry is prohibited. John here with clarity defines idols, and taunts them also, alluding to the words of the prophet in Psalm 114. The idols of the gentiles are silver and gold, the work of men's hands; a mouth they have and speak not, and so forth. Therefore it appears of the matter that images have nothing of religion. For they are of earth, of gold, brass, stone, timber, and so forth. Again, of the form and shape it appears that images are vain. For the form resembles a most gross shape, and even a lie. For neither God nor the saints were of that shape which the idols represent. And now there is no virtue in them. They see not, they hear not, and so forth. How then do they represent God or the saints? I have spoken of idols elsewhere. They that think how there is a diversity between the idols of Christians and those of the gentiles, let them show that theirs are not of wood, or that those other do see, hear, and so forth.

\index{John, Saint}The sins that follow are against the second table, which commands, “You shall not murder, you shall not commit adultery, you shall not steal.” There are many kinds of murders. For they slay most cruelly those who have no sword, but a venomous tongue. Many kill with corrupt doctrine. There are schemers, parricides, and murderers, etc. Unless these repent, they shall not enter into the kingdom of God. And those who swell with envy and malice are homicides, as Saint John said in his canonical Epistle, 1 John 3:1.

Poisoning, witchcraft, or sorcery, or enchanting, pertains to murder. Love potions and enchantments were in the time of Saint John most frequented throughout the Roman Empire; at this day those wicked arts are renewed. But they shall be punished by God, so many as apply themselves to the same.

Fornication also has various forms. This includes whoredom, incest, adultery, and anything else more abominable than these. The Gentiles believed that simple fornication, that is, between two unmarried persons, was no sin. But the Apostle defines the contrary in 1 Corinthians 6:15. This pestilent opinion is revived in many today. Nevertheless, it is certain that a fornicator will not enter the kingdom of God. Ephesians 5.

Finally, this is completed with all its parts. Whereof I spoke once in the exposition of the ten commandments. The Lord Jesus preserve us from all defilement by sin, and so forth. Amen.
